% ---------------------------------------------------------------
\section*{Work, Health and Safety at Mission Systems}
% ---------------------------------------------------------------
 
Mission Systems Pty Ltd is an Australian defence technology company specialising in radar and
sensor systems for uncrewed aerial systems (UAS). The placement involved contributing to the
development and testing of a radar system under a contract with the Australian 
Department of Defence. Work during the placement was split between office-based
engineering at Mission Systems' Lane Cove West office and field trials conducted in
Singleton in the Hunter Valley, NSW. Both environments introduced distinct
health and safety considerations that are addressed in this report.
 
\subsection*{Policies and Regulatory Framework}
 
Mission Systems manages employee health and safety in accordance with the
\textit{Australian Work Health and Safety Act 2011} and relevant state and territory WHS
Regulations, supplemented by the following domain-specific standards and regulatory bodies:
 
\begin{itemize}
    \item \textbf{Australian Radiation Protection and Nuclear Safety Agency (ARPANSA)} ---
    \textit{Radiation Protection Standard for Limiting Exposure to Radiofrequency Fields
    (100\,kHz to 300\,GHz, 2021)}, based on the ICNIRP 2020 recommendations. This standard
    defines mandatory basic restrictions on RF exposure and set the quantitative limits
    governing all radar safety assessments during the placement.
 
    \item \textbf{Australian Communications and Media Authority (ACMA)} --- spectrum licensing
    and authorisation requirements, with Third Party Authorisation (TPA) from the Defence
    Spectrum Office (DSO) required before the radar could be operated outdoors within its
    designated frequency band (15.3--17.2\,GHz).
 
    \item \textbf{Civil Aviation Safety Authority (CASA)} regulations for Remotely Piloted
    Aircraft Systems (RPAS), including Remote Operator Certificate (ReOC) requirements and
    Visual Line of Sight (VLOS) operating conditions for all drone flight activities.
 
    \item \textbf{Mission Systems internal Standard Operating Procedures (SOPs)}, covering
    RPAS flight operations, lithium battery charging and storage, radar payload operation, and
    general office safety.
\end{itemize}
 
The key RF exposure limits relevant to this placement under the ARPANSA standard apply to
the 6--300\,GHz frequency band. For the general public, the maximum time-averaged absorbed
power density is 20\,W/m$^2$ (for averaging intervals of 6 minutes or more), while the
occupational limit is 100\,W/m$^2$. These limits formed the quantitative basis for all RF
safety calculations conducted on the radar-DAA hardware.
 
\subsection*{Training Undertaken}
 
Prior to participating in any radar operations or field activities, the following training
and induction activities were undertaken:
 
\begin{itemize}
    \item \textbf{Office induction} --- familiarisation with the Lane Cove West office
    including the locations of first aid kits, fire extinguishers, designated first aiders,
    fire wardens, and emergency evacuation procedures and assembly points.
 
    \item \textbf{Workstation ergonomics assessment} --- as a significant portion of the
    placement involved sustained computer-based work, a self-assessment of the workstation
    setup was completed in accordance with Mission Systems' ergonomics guidelines, covering
    monitor height, seating posture, keyboard and mouse positioning, and lighting.
 
    \item \textbf{RF and radar safety} --- a briefing covering the relevant ARPANSA exposure
    limits, the methodology for near-field and far-field power density calculations as applied
    to the DAA hardware, and the specific safe-distance requirements derived for both
    the single sub-array and four-sub-array cluster configurations.
 
    \item \textbf{RPAS operations induction} --- familiarisation with Mission Systems' RPAS
    Operations Register and flight SOPs, including pre-flight checklists, emergency procedures,
    the 1:1 altitude-to-distance rule for pilot and ground personnel separation, and NOTAM
    requirements for operations above 400\,ft AGL.
 
    \item \textbf{Lithium battery handling} --- coverage of Mission Systems' SOP for charging,
    storing, and handling Lithium (Li) batteries, including the mandatory use of battery-safe
    enclosures and requirement for F-500 fire extinguishers and fire blankets at all take-off
    and landing sites.
 
    \item \textbf{Physical security training} --- given the Commercial-in-Confidence and
    OFFICIAL classification of project documentation under Contract AFHQ/JDI/14/25, an
    induction into information handling procedures was completed, covering approved data access,
    management of visitors and contractors in the office, and ensuring workstations were clear
    of sensitive information when unattended.
 
    \item \textbf{Cybersecurity awareness} --- training in the identification and reporting of
    phishing attempts, appropriate use of the company VPN when accessing external connections,
    and multi-factor authentication requirements for all company systems.
\end{itemize}
 
% ---------------------------------------------------------------
\section*{Hazard Identification and Risk Management}
% ---------------------------------------------------------------
 
The placement required active engagement with several categories of occupational hazard
across both the office and field environments. The most significant are described below,
alongside the controls applied.
 
\subsection*{Office Ergonomics and Sedentary Work}
 
A substantial portion of the placement involved sustained computer-based work including
software development, data analysis, and documentation. Prolonged sitting without adequate
ergonomic setup can lead to neck, back, and shoulder strain, as well as eye fatigue from
extended screen use.
 
To mitigate these risks, the workstation was configured with a monitor at an appropriate
height to encourage a neutral neck posture, and seating was adjusted to provide adequate
lumbar support. Regular short breaks were taken throughout the working day to reduce eye
strain and encourage movement, including periodic focus changes from screen to distance to
reset natural eye focus. Where possible, working patterns were varied to avoid extended
unbroken sitting.
 
\subsection*{Mental Health and Wellbeing}
 
Periods of intensive project work, particularly in the lead-up to field trials, carry a risk
of overwork and burnout. Awareness of these risks was maintained throughout the placement,
with deliberate effort made to maintain a balanced routine including regular physical activity
outside working hours. Mission Systems' internal culture of open communication between team
members also provided an informal check-in mechanism, consistent with a proactive approach
to workplace mental health.
 
\subsection*{Physical and Cybersecurity}
 
Given that Mission Systems operates under defence contracts with classified deliverables,
physical and information security represented an ongoing WHS-adjacent obligation throughout
the placement. Relevant practices observed included:
 
\begin{itemize}
    \item Ensuring workstations were locked when unattended and that sensitive documents were
    not left visible in shared areas of the office.
    \item Verifying the identity of visitors before granting access to working areas, and
    ensuring external guests were accompanied by staff at all times.
    \item Using multi-factor authentication for all company system logins and connecting via
    VPN when accessing company resources remotely.
    \item Identifying and reporting suspicious emails to the relevant internal contact,
    consistent with training on social engineering and phishing recognition.
\end{itemize}
 
\subsection*{Radiofrequency Radiation Exposure}
 
The primary technical hazard associated with the DAA radar is exposure to
radiofrequency (RF) electromagnetic radiation. Formal safety calculations were conducted and
documented, applying the ARPANSA standard to both the single sub-array and four-sub-array
cluster configurations.
 
For the \textbf{single sub-array} (0.5\,W radiated power), the near-field power density
(within 37\,mm of the antenna face) was calculated to reach up to 1,559\,W/m$^2$ ---
approximately 78 times the general public continuous exposure limit of 20\,W/m$^2$. In the
far field, the minimum safe working distance was determined to be \textbf{220\,mm} from the
radiating face. For the \textbf{four-sub-array cluster} (2\,W total radiated power), the
near-field power density reached up to 1,667\,W/m$^2$ (83 times the public limit), with a
corresponding safe distance of \textbf{775\,mm}.
 
Due to the high operating frequency ($\approx$16.5\,GHz), the skin depth of microwave
radiation in human tissue is between 6 and 12\,mm, limiting any physiological effects to
the body surface. Eye exposure was identified as warranting particular caution, with direct
exposure to the radiating antenna face to be avoided at any range.
 
The following controls were applied:
 
\begin{itemize}
    \item A \textbf{software-based altitude cut-out} was implemented to prevent radar
    activation below a defined altitude threshold during flight, eliminating near-field RF
    exposure to ground personnel during airborne operations by design.
    \item Minimum safe working distances (220\,mm single sub-array; 775\,mm four-sub-array
    cluster) were enforced during all ground handling and bench testing.
    \item Personnel were instructed not to position themselves in the antenna beam path during
    powered testing and not to look directly into the radiating face at any range.
    \item Radar activation during ground tests was announced and controlled, with the test
    area clear of non-essential personnel prior to transmission.
\end{itemize}
 
\subsection*{RPAS Flight Operations}
 
Field testing near Singleton involved the flight of T-drones M1200 uncrewed
aircraft systems (sub-25\,kg RPAS) carrying the radar payload ($\approx$3\,kg), and in some
test serials a second M1200 as an airborne target. The kinetic energy of a sub-25\,kg UAS
in flight represents a significant hazard to personnel and property.
 
Testing was conducted on a private farm property near Singleton, which offered sufficient
open space and low obstructions for the test geometry required. While less remote than a
dedicated military training area, the rural setting still required careful pre-trial planning
to confirm access to power, identify the nearest emergency services, and ensure the test area
was clear of livestock and other personnel not associated with the trial. The following
operational controls were in place:
 
\begin{itemize}
    \item All flights were conducted under Mission Systems' ReOC by accredited remote pilots
    operating within VLOS of their respective aircraft at all times.
    \item The \textbf{1:1 altitude-to-ground-distance rule} was applied for pilot and ground
    personnel positioning throughout all test serials.
    \item Both host and target aircraft were flown on pre-programmed waypoint paths to ensure
    repeatable, ground-truth-referenced trajectories, while pilots retained full manual
    override authority at all times.
    \item A \textbf{500\,m no-fly zone} around the DAA test platform hover position was
    nominated during multi-aircraft operations, with all parallel activities on the base
    formally coordinated and deconflicted prior to flight.
    \item NOTAM filing was completed for operations above 400\,ft AGL, and a DRePL-qualified
    defence pilot was identified as a requirement for any BVLOS operations.
\end{itemize}
 
\subsection*{Lithium Battery Handling}
 
The M1200 platform is powered by a 12S (27\,Ah) Lithium-Ion battery pack at 48\,V,
presenting a fire and chemical hazard if mishandled. Mission Systems' SOP for Charging,
Storing, and Handling Lithium Batteries was followed at all times, including the use of
battery-safe enclosures on the ground and the availability of F-500 fire extinguishers and
fire blankets at both host and target platform take-off and landing sites.
 
% ---------------------------------------------------------------
\section*{Incidents and Preventions During Placement}
% ---------------------------------------------------------------
 
No reportable safety incidents occurred during the placement. The following describes the
principal preventive measures taken across both the office and field environments.
 
\subsection*{Office-Based Work}
 
The most consistent day-to-day safety practice involved maintaining good ergonomic habits
during sustained computer work. On reviewing the initial workstation setup, adjustments were
made to monitor positioning and seating to better align with Mission Systems' ergonomics
guidelines, reducing the risk of developing postural strain over the course of the placement.
Regular breaks and deliberate changes in focus distance were incorporated into the daily
routine to manage eye fatigue. These measures were informed by the initial workstation
assessment and maintained throughout the placement.
 
With respect to cybersecurity, several unsolicited emails were received during the placement
that exhibited characteristics consistent with phishing attempts, including mismatched sender
domains and language designed to create urgency. These were identified using skills developed
during cybersecurity training and were promptly reported to the relevant internal contact,
consistent with company policy.
 
\subsection*{Field Trials at Singleton}
 
The field trial near Singleton was the first opportunity to operate the DAA radar
outdoors at full power within its licensed spectrum band. The most significant preventive
measure was the completion of RF safety calculations prior to departure, ensuring that all
personnel were briefed on safe distances and operational procedures before arriving on site.
 
Pre-trial planning confirmed access to power at the farm property and identified the nearest
emergency services in Singleton township. The test area was inspected on arrival to ensure
it was clear of livestock, overhead powerlines, and other obstructions within the flight
envelope. Battery handling procedures were rehearsed prior to departure, and the required
fire safety equipment was confirmed as part of the pre-departure checklist. No battery
anomalies or other incidents were encountered during the trial.
 
Battery handling procedures were rehearsed prior to the trial, and the required fire safety
equipment was confirmed as part of the pre-departure checklist. No battery anomalies were
encountered during the trial.
 
